\clearpage\fancyhead[L]{\small PALLAS.}\hypertarget{native-374-lower}{}\label{native-374-lower}
\begin{center}\textbf{PALLAS.}\par\vspace{4pt}\hrule\vspace{4pt}\fontsize{10}{12}\selectfont Göttingische gelehrte Anzeigen. No. 176. pp. 1753\;..\;1755. 4 November 1813.\par\vspace{4pt}\hrule\vspace{4pt}\end{center}

\par\indent Last year’s observations gave striking confirmation of the predicted motion of \emph{Pallas}, which was based on the perturbation calculations. We reported on that confirmation at the time in nos. 67 and 127 of this journal [pp. 349 and 356 of this volume]. This confirmation heightened the astronomers’ interest in the planet’s return this year. At our observatory the planet was first seen again on 28 June, exactly at the position given by the ephemeris that Mr. \textsc{Nicolai} had calculated two years earlier. Around the time of opposition. Professor \textsc{Gauss} sought to obtain as many observations as possible, but the exceptionally unfavorable weather last summer yielded very few. Even so, these observations turned out exceptionally well, although the faintness of the planet, again only of the tenth magnitude, meant they were made only with the ring micrometer. Here are all of Professor \textsc{Gauss}’s observations:

\begin{center}\fontsize{10}{12}\selectfont\setlength{\tabcolsep}{5pt}\renewcommand{\arraystretch}{1.15}\begin{tabular}{l|c|c|c}
1813 & Mean time & Right ascension & Northern declination\\\hline
June 28 & $11^{\mathrm u}\;37^{\mathrm m}\;2^{\mathrm s}$ & $327^\circ\;2'\;13''3$ & $14^\circ\;5'\;43''4$\\
Aug. 16 & $10\;36\;23$ & $319\;50\;54.7$ & $.\;.\;.\;.$\\
19 & $9\;31\;0$ & $319\;17\;10.6$ & $10\;0\;8.4$\\
26 & $9\;23\;7$ & $317\;59\;52.0$ & $8\;44\;35.3$\\
\end{tabular}\end{center}

\Needspace{18\baselineskip}\par\indent Mr. \textsc{Burckhardt} and Mr. \textsc{von Lindenau} also sent Professor \textsc{Gauss} the following valuable observations:

\begin{center}Meridian observations of Pallas at the observatory of the Military School in Paris.\end{center}
\begin{center}\fontsize{10}{12}\selectfont\setlength{\tabcolsep}{5pt}\renewcommand{\arraystretch}{1.15}\begin{tabular}{l|c|c|c}
1813 & Sidereal time & Right ascension & Northern declination\\\hline
Aug. 4 & $21^{\mathrm u}\;28^{\mathrm m}\;30^{\mathrm s}5$ & $322^\circ\;7'\;38''0$ & $12^\circ\;11'\;58''5$\\
9 & $21\;24\;42$ & $.\;.\;.\;.$ & $11\;32\;38.4$\\
11 & $.\;.\;.\;.$ & $.\;.\;.\;.$ & $11\;15\;28.1$\\
15 & $21\;20\;7.2$ & $320\;1\;48.0$ & $10\;38\;48.2$\\
17 & $21\;18\;35.2$ & $319\;38\;48.5$ & $10\;19\;29.3$\\
19 & $21\;17\;4.3$ & $319\;16\;4.5$ & $9\;59\;18.3$\\
26 & $21\;11\;56.1$ & $317\;59\;1.8$ & $8\;43\;23.4$\\
\end{tabular}\end{center}

\Needspace{11\baselineskip}\begin{center}\leavevmode\SourceLeaf{375}Observations of Pallas at the Seeberg Observatory:\end{center}

\begin{center}\fontsize{10}{12}\selectfont\setlength{\tabcolsep}{5pt}\renewcommand{\arraystretch}{1.15}\begin{tabular}{l|c|c}
1813 & Mean time & Right ascension\\\hline
Sept. 1 & $10^{\mathrm u}\;25^{\mathrm m}\;16^{\mathrm s}6$ & $316^\circ\;58'\;15''9$\\
3 & $10\;16\;9.6$ & $316\;39\;25.5$\\
4 & $10\;11\;35.4$ & $316\;29\;48.7$\\
\end{tabular}\end{center}

\Needspace{25\baselineskip}\par\indent Professor \textsc{Gauss} passed these observations to Mr. \textsc{Nicolai}. This skilled young astronomer, now employed as Mr. \textsc{von Lindenau}’s assistant at the Seeberg Observatory, first compared them very carefully with the most recently refined elements and found the following close agreement:

\begin{center}Difference from the calculation:\end{center}
\begin{center}\fontsize{10}{12}\selectfont\setlength{\tabcolsep}{7pt}\renewcommand{\arraystretch}{1.15}\begin{tabular}{l|c|c|l}
1813 & \shortstack{Right\\ascension} & Declination & Observer\\\hline
June 28 & $+23''6$ & $+22''4$ & Gauss\\
Aug. 4 & $+36.9$ & $-5.9$ & Burckhardt\\
9 & $.\;.\;.\;.$ & $+5.1$ & Burckhardt\\
11 & $.\;.\;.\;.$ & $+6.2$ & Burckhardt\\
15 & $+26.9$ & $+3.7$ & Burckhardt\\
16 & $+35.6$ & $.\;.\;.\;.$ & Gauss\\
17 & $+28.2$ & $-6.6$ & Burckhardt\\
19 & $+28.5$ & $+4.5$ & Gauss\\
19 & $+24.8$ & $-7.4$ & Burckhardt\\
26 & $+22.4$ & $-4.3$ & Gauss\\
26 & $+18.3$ & $+11.1$ & Burckhardt\\
Sept. 1 & $+22.8$ & $.\;.\;.\;.$ & v. Lindenau\\
3 & $+20.9$ & $.\;.\;.\;.$ & v. Lindenau\\
4 & $+49.9$ & $.\;.\;.\;.$ & v. Lindenau\\
\end{tabular}\end{center}

\Needspace{13\baselineskip}\par\indent From the observations of 15, 17 and 19 August Mr. \textsc{Nicolai} derived the following result for the opposition:

\begin{center}\emph{Ninth observed opposition of Pallas.}\end{center}
\begin{center}1813 August 18. $8^{\mathrm u}\;41^{\mathrm m}\;5^{\mathrm s}$ mean time in Göttingen\end{center}
\begin{center}\fontsize{10}{12}\selectfont\setlength{\tabcolsep}{3pt}\renewcommand{\arraystretch}{1.15}\begin{tabular}{lcr}
True longitude of Pallas & $.\;.\;.\;.\;.\;.\;.\;.$ & $325^\circ\;23'\;56''5$\\
True geocentric latitude & $.\;.\;.\;.\;.\;.\;.$ & $24\;37\;36.1\ \mathrm{N.}$\\
\end{tabular}\end{center}

\noindent Compare this with the prediction in no. 127 of this journal last year [p. 356 of this volume].

\par\indent In these circumstances it did not seem worth the trouble to Professor \textsc{Gauss} to refine the elements again for a still better fit to this opposition.
